\chapter[সিস্টেম নকশা, ডেটা ও প্রিপ্রসেসিং]{সিস্টেম নকশা, ডেটা ও\\ প্রিপ্রসেসিং}\label{ch:system}

\section{কাঠামোর স্থাপত্য}
কাঠামোটি এমন স্তরে সাজানো যা প্রস্তাবনার কর্মপ্রবাহের প্রতিফলন (\cref{fig:architecture})। প্রতিটি স্তর একটি Python মডিউল যার একটিই
দায়িত্ব। মডিউলগুলো কেবল স্থিতিশীল ফরম্যাটের ফাইলের (Parquet, SQLite ও JSON) মাধ্যমে যোগাযোগ করে, তাই অন্যগুলো না বদলে যেকোনো
স্তর প্রতিস্থাপন করা যায়। যেমন, একই SQLite ডেটাবেস পড়ে Streamlit ড্যাশবোর্ডের বদলে Power BI ব্যবহার করা যায়।

\begin{figure}[htbp]\centering
\begin{tikzpicture}[node distance=3.5mm, every node/.style={font=\footnotesize},
  L/.style={draw=brand, rounded corners, fill=brand!7, minimum width=12.5cm, minimum height=8.5mm, align=left, text width=12cm},
  arr/.style={-{Stealth[length=2mm]}, thick, brand}]
\node[L] (a) {\textbf{১. ডেটা উৎস} \hfill কাঁচা CSV এক্সপোর্ট (\NRows লেনদেন, \NCols কলাম)};
\node[L, below=of a] (b) {\textbf{২. গ্রহণ ও যাচাই} \hfill টাইপসহ লোড, স্কিমা, নিয়ম ও অভেদ যাচাই \hfill \texttt{data.py}};
\node[L, below=of b] (c) {\textbf{৩. পরিষ্কারকরণ ও সংরক্ষণ} \hfill imputation, ডুপ্লিকেট অপসারণ, outlier চিহ্ন, স্টার স্কিমা \hfill \texttt{data.py}};
\node[L, below=of c] (d) {\textbf{৪. অনুসন্ধানমূলক অ্যানালিটিক্স} \hfill প্রবণতা, মৌসুমিতা, টেস্ট, elasticity \hfill \texttt{eda.py}};
\node[L, below=of d] (e) {\textbf{৫. ফিচার ইঞ্জিনিয়ারিং} \hfill leakage-মুক্ত lag, rolling পরিসংখ্যান, ক্যালেন্ডার, মূল্য \hfill \texttt{features.py}};
\node[L, below=of e] (f) {\textbf{৬. মেশিন লার্নিং} \hfill পূর্বাভাস $\cdot$ ক্লাসিফিকেশন $\cdot$ ক্লাস্টারিং $\cdot$ অ্যানোমালি ডিটেকশন};
\node[L, below=of f] (g) {\textbf{৭. মূল্যায়ন ও ব্যাখ্যা} \hfill বেঞ্চমার্ক, মেট্রিক, permutation importance, SHAP};
\node[L, below=of g] (h) {\textbf{৮. উপস্থাপন} \hfill Streamlit ড্যাশবোর্ড $\cdot$ প্রকল্প ওয়েবসাইট $\cdot$ \LaTeX{} প্রতিবেদন \hfill \texttt{dashboard/}};
\node[L, below=of h, fill=accent!10, draw=accent] (i) {\textbf{৯. সিদ্ধান্ত-সহায়তা} \hfill স্বয়ংক্রিয় অন্তর্দৃষ্টি ও সুপারিশ \hfill \texttt{insights.py}};
\foreach \x/\y in {a/b,b/c,c/d,d/e,e/f,f/g,g/h,h/i} \draw[arr] (\x)--(\y);
\end{tikzpicture}
\caption{বুদ্ধিমান অ্যানালিটিক্স কাঠামোর স্তরভিত্তিক স্থাপত্য}\label{fig:architecture}
\end{figure}

একটি কমান্ড, \texttt{python run\_pipeline.py}, সব স্তর ক্রমানুসারে চালায় এবং লেখে:
\begin{itemize}
  \item পরিষ্কার করা ডেটাসেট (\path{data/processed}) ও ওয়্যারহাউস (\path{data/warehouse/dancake.db});
  \item JSON-এ মেট্রিক (\path{outputs/metrics}) ও CSV-তে সারণি (\path{outputs/tables});
  \item এই প্রতিবেদনের প্রতিটি চিত্র ও \LaTeX{} সারণি (\path{report/figures}, \path{report/tables}), এবং লেখায় উদ্ধৃত প্রতিটি
        সংখ্যা ধারণকারী \LaTeX{} macro-র একটি ফাইল (\path{report/generated/numbers.tex});
  \item ড্যাশবোর্ড ও ওয়েবসাইটের জন্য সংক্ষিপ্ত সমষ্টিগত ডেটা (\path{data/dashboard})।
\end{itemize}
কোনো ফলাফল হাতে টাইপ করা হয় না, তাই প্রতিবেদন, ড্যাশবোর্ড ও ওয়েবসাইটের মধ্যে অমিল হতে পারে না। নতুন মাসের ডেটা যোগ করতে
কেবল একটি কমান্ড আবার চালাতে হয়।

\section{ডেটাসেট}
\Cref{tab:dqsummary} ডেটাসেট ও যাচাই-স্তরের ফলাফলের সারসংক্ষেপ দেয়। প্রতিটি লেনদেন 13টি ক্যাটাগরির \NProducts টি পণ্যের একটির,
যা \NRegions টি অঞ্চলে \NChannels টি চ্যানেলের (General Trade, Modern Trade, Distributor, E-commerce) মাধ্যমে \NCustTypes ধরনের
গ্রাহকের (Retailer, Super Shop, Distributor, Corporate, Online Customer) কাছে বিক্রি হয়। পাঁচটি পণ্য 2020-এর পরে চালু হয়েছে:
লাচ্ছা সেমাই, দুটি পাউরুটি, একটি ক্রোয়াসাঁ এবং একটি ক্যাপুচিনো মাফিন গিফট প্যাক।

\begin{table}[htbp]\centering\small
\caption{ডেটাসেটের পরিচিতি ও ডেটা-মান যাচাইয়ের ফলাফল}\label{tab:dqsummary}
\begin{tabular}{@{}lr@{}}
\toprule
\textbf{যাচাই} & \textbf{ফলাফল} \\ \midrule
সারি $\times$ কলাম & \NRows $\times$ \NCols \\
আচ্ছাদিত সময়কাল & জানুয়ারি 2020 -- সেপ্টেম্বর 2026 \\
পণ্য / অঞ্চল / চ্যানেল / গ্রাহকের ধরন & \NProducts / \NRegions / \NChannels / \NCustTypes \\
অনুপস্থিত সেল (``NA'', ``-'' জাতীয় লুকানো মানসহ) & 0 \\
হুবহু ডুপ্লিকেট সারি / ডুপ্লিকেট লেনদেন আইডি & 0 / 0 \\
ডোমেইন-নিয়ম লঙ্ঘন (পরিমাণ, মূল্য, ডিসকাউন্ট, মূল্যমান) & 0 \\
মোট $=$ পরিমাণ $\times$ বাজার মূল্য (সর্বোচ্চ পরম ত্রুটি) & $<10^{-11}$ টাকা \\
নিট $=$ মোট $-$ ডিসকাউন্ট (সর্বোচ্চ পরম ত্রুটি) & $<10^{-11}$ টাকা \\
তারিখের সঙ্গে বছর/মাস কলামের সামঞ্জস্য & 100\% \\
পণ্য আইডি $\rightarrow$ একটি নাম, ক্যাটাগরি, প্যাক সাইজ & 100\% \\
পরিসংখ্যানিক পরিমাণ-outlier (চিহ্নিত, রেখে দেওয়া) & \NOutlierFlags (\OutlierPct\%) \\
ওয়্যারহাউস মিলকরণ (SQL মোট $-$ pandas মোট) & 0.00 টাকা \\ \bottomrule
\end{tabular}
\end{table}

\section{প্রিপ্রসেসিং পাইপলাইন}
পাইপলাইনটি রক্ষণাত্মকভাবে লেখা: এটিকে কেবল সরবরাহকৃত (ব্যতিক্রমীভাবে পরিষ্কার) ফাইলে নয়, অপরিষ্কার ডেটাতেও সঠিকভাবে কাজ করতে
হবে। অপরিষ্কার ডেটায় এর আচরণ দুইভাবে যাচাই করা হয়েছে। ছয়টি হাতে-কলমে অনুশীলন স্ক্রিপ্ট ইচ্ছাকৃতভাবে ডেটার অনুলিপি নষ্ট করে
মেপেছে কতটা মেরামত হয়। এছাড়া \texttt{tests/}-এর স্বয়ংক্রিয় টেস্টগুলো অনুপস্থিত মান, ডুপ্লিকেট ও অসম্ভব মান ঢুকিয়ে নিশ্চিত করে যে
পাইপলাইন সেগুলো মেরামত করে (\cref{app:reproduce})। ধাপগুলো একটি নির্দিষ্ট ক্রমে প্রয়োগ হয়, কারণ কিছু ধাপ অন্যগুলোর উপর নির্ভরশীল।

\subsection{ডেটা টাইপ রূপান্তর ও যাচাই}
তারিখ একটি সুনির্দিষ্ট ফরম্যাটে (\texttt{dd-mm-yy}) পার্স করা হয় এবং সংখ্যাসূচক কলাম রূপান্তর করা হয়; যেসব মান পার্স হয় না সেগুলো
অনুপস্থিত হয়ে যায়, ফলে কখনো নীরবে ভুল পড়া হয় না। টেক্সট ফিল্ডের অতিরিক্ত স্পেস মুছে ফেলা হয়। এরপর যাচাই ধাপ একটি
যন্ত্র-পাঠযোগ্য মান-প্রতিবেদন তৈরি করে (\path{outputs/metrics/data_quality.json})। এটি অনুপস্থিত মান, লুকানো placeholder,
ডুপ্লিকেট সারি ও কী, ডোমেইন-নিয়ম লঙ্ঘন, কলামগুলোর মধ্যে হিসাবগত অভেদ, এবং ক্যালেন্ডার ফিল্ড ও পণ্য-বৈশিষ্ট্যের সামঞ্জস্য যাচাই করে।

\subsection{অনুপস্থিত মান ব্যবস্থাপনা}
\citet{garcia2015} অনুসরণ করে একটি সিদ্ধান্ত-স্তরক্রম বাস্তবায়িত হয়েছে:
\begin{enumerate}
  \item যেখানে অভেদ আছে সেখানে অন্য কলাম থেকে মানটি হুবহু \emph{নির্ণয়} করা (যেমন
        $\text{পরিমাণ} = \text{মোট}/\text{বাজার মূল্য}$, বা তারিখ থেকে বছর ও মাস);
  \item কোনো অপরিহার্য শনাক্তকারী (লেনদেন আইডি, তারিখ, পণ্য) অনুপস্থিত হলে সারিটি \emph{বাদ} দেওয়া, কারণ অনুমান করা হবে বানানো তথ্য;
  \item সংখ্যাসূচক ফাঁক একই পণ্যের median দিয়ে \emph{পূরণ} করা, যা বিষম বণ্টনে robust;
  \item পণ্য-বৈশিষ্ট্য পণ্যের mode দিয়ে পূরণ করা, এবং অনুপস্থিত বাজার-বৈশিষ্ট্য স্পষ্টভাবে ``Unknown'' লেবেল দেওয়া যাতে তা দৃশ্যমান থাকে।
\end{enumerate}
এলোমেলোভাবে 3\% সেল মুছে ফেলা অনুশীলন-পরীক্ষায় হিসাবগত অভেদ থেকে নির্ণয় পরিমাণ, মূল্য ও বিক্রয়মূল্য হুবহু ফিরিয়ে এনেছে, অথচ সরল
mean imputation প্রতি অর্ডারে কয়েক ইউনিট ত্রুটি তৈরি করেছে।

\subsection{ডুপ্লিকেট অপসারণ}
প্রথমে হুবহু ডুপ্লিকেট সরানো হয়। এরপর ডুপ্লিকেট কী (একই লেনদেন আইডি কিন্তু ভিন্ন মান) সবচেয়ে সম্পূর্ণ সংস্করণটি রেখে সমাধান করা
হয়। আইডি বড় হাতের অক্ষরে রূপান্তর ও স্পেস মুছে তুলনা করা হয়, যাতে `` dc2020000001''-এর মতো লুকানো ডুপ্লিকেট ধরা পড়ে।
প্রায়-ডুপ্লিকেট (একই তারিখ, পণ্য, অঞ্চল, চ্যানেল ও পরিমাণ কিন্তু ভিন্ন আইডি) কেবল \emph{চিহ্নিত} করা হয়, কারণ বিক্রয় ডেটায় এগুলো
সাধারণত প্রকৃত পুনরাবৃত্ত অর্ডার।

\subsection{Outlier শনাক্তকরণ ও ব্যবস্থাপনা}
অর্ডারের পরিমাণ তীব্রভাবে ডানদিকে বিষম (skewness \QtySkew; median \QtyMedian ইউনিট), এবং প্রতিটি পণ্যের নিজস্ব স্বাভাবিক অর্ডার-আকার
আছে। তাই একটি সর্বজনীন $z$-score ছোট পণ্যের অনেক বেশি এবং বড় পণ্যের অনেক কম অর্ডার চিহ্নিত করে। অনুশীলন-পর্যায়ে চারটি পদ্ধতির
তুলনা করা হয়েছে: IQR নিয়ম, $z$-score, MAD-ভিত্তিক modified $z$-score, এবং log স্কেলে প্রতি-পণ্য IQR নিয়ম। শেষটিই সবচেয়ে
ভারসাম্যপূর্ণ, এবং পাইপলাইনে এটিই ব্যবহৃত:
\[
  \text{flag}_i = \bigl[\, \log(1+q_i) \notin [\,Q_1^{(p)} - 1.5\,\text{IQR}^{(p)},\; Q_3^{(p)} + 1.5\,\text{IQR}^{(p)}\,] \bigr],
\]
যেখানে প্রতিটি পণ্য $p$-এর জন্য চতুর্থক আলাদাভাবে হিসাব করা হয়। অসম্ভব মান (যেমন এক-এর কম পরিমাণ বা 50\%-এর বেশি ডিসকাউন্ট)
ত্রুটি হিসেবে গণ্য: মুছে আবার নির্ণয় করা হয়। অস্বাভাবিক কিন্তু সম্ভব মান চিহ্নিত করে রাখা হয়। প্রকৃত ডেটায় এমন \NOutlierFlags টি
চিহ্ন পাওয়া গেছে (অর্ডারের \OutlierPct\%), সবই বিশ্বাসযোগ্য বাল্ক অর্ডার, এবং কোনো অসম্ভব মান পাওয়া যায়নি।

\subsection{Normalisation ও এনকোডিং}
Scaling সর্বজনীনভাবে নয়, বরং প্রতিটি মডেলের পাইপলাইনের \emph{ভেতরে} কেবল প্রশিক্ষণ-সারিতে fit করে প্রয়োগ করা হয়, যাতে leakage না
হয়। গণনাজাতীয় ফিচার $\log(1+x)$ দিয়ে রূপান্তরিত হয় ভেদাঙ্ক স্থিতিশীল করতে, তারপর স্কেল-সংবেদনশীল মডেলের (linear regression,
KNN ও SVR) জন্য standardise করা হয়। নামবাচক ক্যাটাগরি (অঞ্চল, চ্যানেল, গ্রাহকের ধরন, ক্যাটাগরি) one-hot এনকোড করা হয়। মাসকে
চক্রাকারে $(\sin 2\pi m/12, \cos 2\pi m/12)$ হিসেবে এনকোড করা হয়, যাতে ডিসেম্বর ও জানুয়ারি প্রতিবেশী হয়। উচ্চ-সংখ্যক মানের
পারস্পরিক ক্রিয়া (চ্যানেল $\times$ অঞ্চল, পণ্য) কেবল প্রশিক্ষণ ডেটা দিয়ে additive smoothing-সহ target-encode করা হয়,
\[
  \tilde\mu_k = \frac{n_k \bar y_k + m\,\bar y}{n_k + m}, \qquad m=50,
\]
ফলে বিরল মানগুলো সামগ্রিক গড়ের দিকে সংকুচিত হয়।

\subsection{চাহিদা পূর্বাভাসের জন্য ফিচার ইঞ্জিনিয়ারিং}
লেনদেনগুলোকে \PanelSeries টি পণ্য $\times$ মাসের একটি \emph{প্যানেলে} একত্র করা হয়েছে। মাস $t$-এর প্রতিটি ফিচার কেবল মাস $t-1$-এর
শেষে পাওয়া তথ্য, অথবা আগে থেকে পরিকল্পিত তথ্য ব্যবহার করে (\cref{tab:features})। প্রশিক্ষণ, পরীক্ষা ও ভবিষ্যৎ পূর্বাভাসে একই ফিচার
ফাংশন ব্যবহৃত, তাই নিয়মটি এক জায়গায় কার্যকর হয়। একটি শিক্ষণীয় যাচাইয়ে সঠিক তিন-মাসের rolling গড়কে (rolling-এর আগে shift)
বর্তমান মাস অন্তর্ভুক্ত একটি ``leaky'' সংস্করণের সঙ্গে তুলনা করা হয়েছে। Leaky সংস্করণের সঙ্গে লক্ষ্য-চলকের সম্পর্ক বেশি শক্তিশালী ছিল,
আর ঠিক এ কারণেই leakage প্রলুব্ধকর ও বিপজ্জনক \citep{kaufman2012}।

\begin{table}[htbp]\centering\small
\caption{চাহিদা পূর্বাভাস মডেলের ফিচার (মোট \PanelFeatures টি)}\label{tab:features}
\begin{tabularx}{\linewidth}{@{}lX@{}}
\toprule
\textbf{দল} & \textbf{ফিচার} \\ \midrule
ইতিহাস (lag) & মাস $t-1, t-2, t-3, t-6, t-12$-এর ইউনিট \\
Rolling পরিসংখ্যান & আগের 3, 6 ও 12 মাসের গড়; আগের 3 মাসের standard deviation \\
গতিবেগ & গত মাসের বছর-ভিত্তিক প্রবৃদ্ধি; 12-মাসের প্রবণতা-অনুপাত \\
মূল্য (পরিকল্পিত) & মাস $t$-এর তালিকা মূল্য; বছর-ভিত্তিক তালিকা-মূল্য পরিবর্তন \\
পণ্য & ক্যাটাগরি (one-hot); গ্রামে প্যাকের ওজন; গিফট-প্যাক সূচক; চালুর পর থেকে মাস \\
ক্যালেন্ডার & মাসের নম্বর; চক্রাকার মাস (sine, cosine); ডিসেম্বর ও ফেব্রুয়ারি সূচক; মাসের দিনসংখ্যা; সময়-সূচক \\ \bottomrule
\end{tabularx}
\end{table}

\section{ডাইমেনশনাল ডেটাবেস}
পরিষ্কার করা ডেটা SQLite-এ একটি স্টার স্কিমায় লোড করা হয় (\cref{fig:star})। Fact table \texttt{fact\_sales}-এ প্রতিটি লেনদেনের জন্য
এর পরিমাপসহ একটি সারি থাকে। Dimension table-গুলো পণ্য, তারিখ (বাংলাদেশের মতো শুক্র--শনি সাপ্তাহিক ছুটির চিহ্নসহ) এবং অঞ্চল,
চ্যানেল ও গ্রাহকের ধরন মিলিয়ে গঠিত 220টি বাজার-বিভাগ বর্ণনা করে। Foreign key-তে index থাকায় সাধারণ KPI কোয়েরি এক সেকেন্ডেরও কম
সময়ে ফল দেয়। একটি মিলকরণ যাচাই নিশ্চিত করে যে ওয়্যারহাউসের মোট নিট বিক্রয় পরিষ্কার ডেটাসেটের মোটের সঙ্গে পয়সা পর্যন্ত মেলে।
\Cref{app:sql}-এ উদাহরণ বিশ্লেষণী SQL কোয়েরি দেওয়া আছে।

\begin{figure}[htbp]\centering
\begin{tikzpicture}[every node/.style={font=\scriptsize},
  T/.style={draw=brand, rectangle split, rectangle split parts=2, rounded corners, fill=white, text width=3.4cm, align=left},
  F/.style={T, draw=accent, fill=accent!6, text width=3.8cm}]
\node[F] (f) at (0,0) {\textbf{fact\_sales} (\NRows)\nodepart{second}transaction\_id\\date\_id (FK)\\product\_id (FK)\\segment\_id (FK)\\list/market/unit price\\discount\_pct, quantity\\gross, discount, net sales\\flag\_outlier\_qty};
\node[T] (p) at (-5.2,1.6) {\textbf{dim\_product} (\NProducts)\nodepart{second}product\_id (PK)\\product\_name, category\\pack\_size, launch\_year\\mrp\_bdt};
\node[T] (d) at (5.2,1.6) {\textbf{dim\_date} (2,465)\nodepart{second}date\_id (PK), date\\year, quarter, month\\day\_of\_week, is\_weekend};
\node[T] (s) at (0,-4.1) {\textbf{dim\_segment} (220)\nodepart{second}segment\_id (PK)\\region, channel\\customer\_type};
\draw[thick, brand] (f.west) -- (p.east); \draw[thick, brand] (f.east) -- (d.west); \draw[thick, brand] (f.south) -- (s.north);
\end{tikzpicture}
\caption{বিশ্লেষণী ডেটাবেসের স্টার স্কিমা}\label{fig:star}
\end{figure}
